% Slide "Several cores: the shared memory becomes the limit": MLUP/s of all cores together, lid-driven cavity (SweepgenLdcCPU), one 64^3 block
% per MPI process, best of 3 runs of 200 steps. 1, 4 and 8 cores from one session (board,
% scripts/presentation4/all.sh, logs docs/presentation4/board-logs/board/
% ldc-{gcc,clang}-{scalar,autovec,rvv}-{f32,f64}-{1,4,8}core.log), printed by
% python3 scripts/presentation4/slide_numbers.py --tex docs/presentation4/board-logs/board
% Bars and labels: GCC 14.2. Diamonds: Clang 18.1.
% plain scalar = scalar, best scalar = scalar_cse_aa, auto-vec. in place = autovec_cse_aa
% (hatched), RVV in place + prefetch = explicit_cse_aa_pf, RVV split loop = explicit_cse_aa_split (in place,
% no software prefetch: 0 hints).
% Memory traffic of RVV in place + prefetch (184 bytes per cell in f32, 368 in f64, the AA slide's
% numbers; walls not counted), GCC: 4 cores 17.91*184 = 3.3 / 9.89*368 = 3.6 GB/s, 8 cores
% 19.40*184 = 3.6 / 10.29*368 = 3.8 GB/s. (The split loop passes over the PDFs more than once:
% its traffic per cell is not 184 bytes, so no GB/s is given for it.)
% Both panels share ymax=27 so f32 and f64 bars compare by height. The speedups of the same runs:
% figures/speedup-multicore-{f32,f64}.tex.
\begin{tikzpicture}
  %   Styles mc, pl, bs, av, rv, sp, cl, cllegend: preamble of slides.tex
  \begin{axis}[mc, ymax=27, name=a32, title={f32}, ylabel={MLUP/s, all cores together}]
    \addplot[pl] coordinates {(1,1.38)[1.4] (4,5.27)[5.3] (8,8.77)[8.8]};
    \addplot[bs] coordinates {(1,2.00)[2.0] (4,7.30)[7.3] (8,12.50)[12.5]};
    \addplot[av] coordinates {(1,2.01)[2.0] (4,7.27)[7.3] (8,12.52)[12.5]};
    \addplot[rv] coordinates {(1,6.56)[6.6] (4,17.91)[17.9] (8,19.40)[19.4]};
    \addplot[sp] coordinates {(1,7.68)[7.7] (4,18.31)[18.3] (8,19.49)[19.5]};
    \addplot[cl,xshift=-12pt] coordinates {(1,1.36) (4,5.23) (8,8.54)};
    \addplot[cl,xshift=-6pt] coordinates {(1,2.11) (4,7.66) (8,12.91)};
    \addplot[cl,xshift=0pt] coordinates {(1,2.10) (4,7.66) (8,12.90)};
    \addplot[cl,xshift=6pt] coordinates {(1,6.19) (4,17.19) (8,18.79)};
    \addplot[cl,xshift=12pt] coordinates {(1,7.54) (4,17.90) (8,19.21)};
  \end{axis}
  %   The legend sits in the f64 panel: its bars leave the top half free
  \begin{axis}[mc, ymax=27, at={(a32.east)}, anchor=west, xshift=1.0cm, title={f64},
      legend style={font=\tiny,at={(0.98,0.98)},anchor=north east,draw=psgray!50},
      legend columns=2]
    \addplot[pl] coordinates {(1,0.92)[0.9] (4,3.39)[3.4] (8,5.17)[5.2]};
    \addplot[bs] coordinates {(1,1.46)[1.5] (4,5.16)[5.2] (8,7.59)[7.6]};
    \addplot[av] coordinates {(1,1.44)[1.4] (4,5.16)[5.2] (8,7.63)[7.6]};
    \addplot[rv] coordinates {(1,3.89)[3.9] (4,9.89)[9.9] (8,10.29)[10.3]};
    \addplot[sp] coordinates {(1,4.71)[4.7] (4,10.06)[10.1] (8,10.35)[10.4]};
    \addplot[cl,xshift=-12pt] coordinates {(1,0.89) (4,3.39) (8,5.15)};
    \addplot[cl,xshift=-6pt] coordinates {(1,1.58) (4,5.49) (8,7.79)};
    \addplot[cl,xshift=0pt] coordinates {(1,1.58) (4,5.59) (8,7.86)};
    \addplot[cl,xshift=6pt] coordinates {(1,3.63) (4,9.43) (8,9.95)};
    \addplot[cl,xshift=12pt,forget plot=false,cllegend] coordinates {(1,4.43) (4,9.79) (8,10.18)};
    \legend{plain scalar, best scalar, auto-vec.\ in place, RVV in place + prefetch, RVV split loop,
      Clang}
  \end{axis}
\end{tikzpicture}
